\templateentry
{Fenwick Tree}
{Soma prefixos e faz update pontual em $O(\log n)$.}
{Atua com um array de inteiros $a_0, a_1, \dots, a_{n - 1}$. Suporta dois tipos de operação:
\begin{itemize}
	\item $add(i, val)$: adiciona o valor $val$ ao índice $i$ ($a_i \to a_i + val$). Complexidade: $O(\log(n))$.
	\item $sum(L, R)$: Calcula a soma $a_L + a_{L + 1} + \cdots + a_{R - 1}$. Complexidade: $O(\log(n))$.
\end{itemize}}
{dsa/fenwick-tree.cpp}

\templateentry
{Segtree}
{Monoide genérico com update pontual e query em intervalo.}
{Atua com uma operação $\oplus$ em elementos de uma monoide $S$ com elemento neutro $e$, em um array $a_1, a_2, \dots, a_{n - 1}$ com $a_i \in S$. Suporta dois tipos de operações:
\begin{itemize}
	\item $set(i, val)$: modifica o array na posição $i$ ($a_i \to val$). Complexidade: $O(\log(n))$.
	\item $sum(L, R)$: calcula $a_l \oplus a_{l + 1} \oplus \dots \oplus a_{r - 1}$. Complexidade: $O(\log(n))$.
\end{itemize}}
{dsa/segtree.cpp}

\templateentry
{Lazy Segtree}
{Segtree com lazy propagation para updates em intervalo.}
{Similar a uma segtree normal, mas suporta modificações em range. Suporta três tipos de operações:
\begin{itemize}
	\item $set(i, val)$: modifica o array na posição $i$ ($a_i \to val$). Complexidade: $O(\log(n))$.
	\item $sum(L, R)$: calcula $a_l \oplus a_{l + 1} \oplus \dots \oplus a_{r - 1}$. Complexidade: $O(\log(n))$.
	\item $apply(L, R, U)$: Aplica a função $mapping$ no range $[L, R)$ ($a_i \to mapping(U, a_i)$, para $i = L, L + 1, \dots, R - 1$). Complexidade: $O(\log(n))$.
\end{itemize}}
{dsa/lazy-segtree.cpp}

\templateentry
{Sparse Table}
{Consultas estáticas idempotentes em $O(1)$ após preprocessamento.}
{Permite calcular soma de funções idempotentes em um intervalo de um array. Ex: Calcular o mínimo no intervalo $[L, R)$
Complexidade:
\begin{itemize}
	\item Construção: $O(n\log(n))$.
	\item Consulta: $O(1)$.
\end{itemize}
Outros exemplos de funções incluem $gcd$, $lcm$ e $max$.}
{dsa/sparse-table.cpp}

\templateentry
{Disjoint Set Union}
{Merge/find amortizado quase constante para componentes.}
{Representa uma lista de conjuntos, e suporta dois tipos de operações:
\begin{itemize}
	\item $merge(X, Y)$: Junta os conjuntos dos números $X$ e $Y$.
	\item $root(X)$: encontra o elemento representante do conjunto de $X$. Dois números estão no mesmo conjunto quando eles tem o mesmo representante.
\end{itemize}}
{dsa/dsu.cpp}

\subsection{Convolução}
\varianttext
{FFT e NTT para convolução em $O(n\log n)$.}
{Calcula a convolução de dois arrays em $O(n\log(n))$.
A convolução de dois arrays $a$ e $b$ é definida como um array $c$ tal que:
\[c_k = \sum\limits_{i + j = k} a_i b_j\]}
\subsubsection{FFT}
\addtemplate{dsa/fft.cpp}
\subsubsection{NTT}
\addtemplate{dsa/ntt.cpp}

\templateentry
{Max Add Segtree}
{}
{Segtree que permite soma em range, e calcular o máximo em um range}
{dsa/max-add-segtree.cpp}

\templateentry
{Zeta e Mobius}
{}
{Zeta: dada uma função $f: 2^{[1..n]} \to S$, encontra uma função $g: 2^{[1..n]} \to S$ tal que \[g(A) = \sum\limits_{B \subseteq A} f(B)\]
Mobius: realiza a transformação inversa.}
{dsa/zeta.cpp}

\templateentry
{Árvore de Li Chao}
{}
{Suporta dois tipos de operações:
\begin{itemize}
	\item $insert(line)$: adiciona a reta $line (y = ax + b)$ ao conjunto de retas. Complexidade: $O(\log(N))$.
	\item $query(X)$: encontra a reta do conjunto que minimiza o valor $aX + b$. Complexidade: $O(\log(N))$.
\end{itemize}
Onde $N$ é o tamanho do intervalo em que as retas atuam.}
{dsa/li-chao.cpp}

\templateentry
{Árvore de Li Chao++}
{}
{Suporta dois tipos de operações:
\begin{itemize}
	\item $insert(line, L, R)$: adiciona a reta $line (y = ax + b)$ ao conjunto de retas, atuando no intervalo $L, R$. Complexidade: $O(\log^2(R - L))$.
	\item $query(X)$: encontra a reta do conjunto que minimiza o valor $aX + b$. Complexidade: $O(\log(L))$.
\end{itemize}}
{dsa/li-chao-extended.cpp}

\templateentry
{Árvore Cartesiana}
{}
{Constrói uma árvore cartesiana dos inteiros $0, 1, \dots, n - 1$ a partir de um critério de comparação.
A árvore é retornada no formato de um array $par$ de tal forma que $par_i$ é o pai do nó $i$ na árvore.
A raiz da árvore é o nó mínimo de acordo com o critério de comparação.
Complexidade: $O(n)$.}
{dsa/cartesian-tree.cpp}

\templateentry
{Árvore de Busca Binária Aleatória}
{}
{}
{dsa/rbst.cpp}

\templateentry
{Treap implícita}
{}
{Treap implícita para sequência: suporta \texttt{split}/\texttt{merge} por posição, inserção, remoção, reversão lazy e soma em intervalo.}
{dsa/treap.cpp}

\templateentry
{Treap de conjunto ordenado}
{}
{Treap por chave para conjunto ordenado: inserção, remoção, busca, $k$-ésimo elemento, contagem $\le x$, predecessor e sucessor.}
{dsa/treap-set.cpp}
